% Source: https://artofproblemsolving.com/community/c3h3607860_simple_trig_problem
% Stable identity: aops-topic-3607860-q01
Every solution has $|m|\le100$, and the equation is odd, so nonzero solutions occur in opposite pairs. For $m>0$ only positive sine lobes can contain roots. The lobe $(0,\pi)$ contributes one positive root: $f(m)=\sin m-\dfrac{m}{100}$ starts at zero with positive derivative, is strictly concave inside the lobe, and ends negative. For each lobe $(2k\pi,(2k+1)\pi)$ with $k=1,\ldots,15$, $f$ is negative at both ends and positive at its midpoint, since that midpoint is at most $\dfrac{61\pi}{2}<100$. Strict concavity and the derivative $f'(m)=\cos m-\dfrac{1}{100}$, which decreases across each such lobe, imply exactly two crossings. The next positive lobe starts at $32\pi>100$. Thus there are $1+2(15)=31$ positive roots, $31$ negative roots and the root $0$, for $\boxed{63}$ in all.
